% Source: https://artofproblemsolving.com/community/c3h2116882_some_sequences
% Stable identity: aops-topic-2116882-q03
\textbf{(a)} For $n=2$ we need $2^k+2=34$, so $2^k=32$ and $k=5$. The rule $n^5+n$ fits all four terms: $1+1=2$, $32+2=34$, $243+3=246$ and $1024+4=1028$. The fifth term is
\[5^5+5=3125+5=\boxed{3130}.\]

\textbf{(b)} A polynomial of degree at most $3$ has constant third differences. The differences of $2,34,246,1028$ are
\[32,\ 212,\ 782;\qquad\text{second differences } 180,\ 570;\qquad\text{third difference } 390 .\]
Keeping the third difference $390$, the next second difference is $570+390=960$, the next difference is $782+960=1742$, and the fifth term is
\[1028+1742=\boxed{2770}.\]

So both rules fit the four given terms but they give different fifth terms. A ``next term'' is only determined once a rule has been fixed; $n^5+n$ is the simpler and intended pattern.
